\section*{Introduction}

Nutrient profiling systems rank foods according to population-level nutritional
criteria and support front-of-package labelling, dietary guidance, policy design
and food-system monitoring. Their public-health value depends on stability:
foods that align with established dietary guidance should remain recognisable as
healthier choices at the population level. Yet this same stability limits their
use in precision nutrition, where individuals differ in postprandial glucose,
lipid and inflammatory responses to the same foods
\cite{zeevi2015personalized,berry2020postprandial}. A precision-ready nutrient
profiling system would benefit from a structure that preserves the population
nutrition prior while allowing bounded individual calibration.

We propose personalized calibration as a general transformation of nutrient
profiling systems. Rather than replacing a static score with an unconstrained
recommendation model, the framework anchors each food to a population-level NPS
prior and adds an individual-specific deviation within a pre-specified range.
This design asks whether a static food score can remain nutritionally coherent
at the population mean while expressing biologically interpretable heterogeneity
across individuals.

GMNPS implements this design with gut microbiome-informed nutrient responses as
a proof-of-concept calibration layer. The gut microbiome is a plausible
calibration signal because it links habitual diet, microbial metabolism and host
metabolic traits in deeply phenotyped cohorts
\cite{asnicar2021microbiome,asnicar2025gut}. We do not assume that the
microbiome is the only possible personalization layer. Instead, we use it to
test whether an interpretable biological layer can move NPS toward precision
nutrition without weakening population-level dietary guidance.

Food Compass 2.0 was selected as the baseline prior because it provides broad
and current NPS attribute coverage across foods and beverages
\cite{mozaffarian2021foodcompass,barrett2024foodcompass2}. In the GMNPS
framework, this baseline is not the main contribution. It supplies the static
population score around which personalized calibration is bounded. The central
question is whether the calibrated score preserves this prior while revealing
individual variation that a static NPS cannot represent.

\section*{Methods}

\subsection*{Study design}

GMNPS was designed as an anchored personalized calibration framework for
nutrient profiling. For individual \(i\) and food \(j\), the score is
\[
  \mathrm{GMNPS}_{ij} = \mathrm{clip}_{1,100}\left(\mathrm{NPS}_{j} + D_{ij}\right),
\]
where \(\mathrm{NPS}_{j}\) is the baseline population score and \(D_{ij}\) is a
bounded personalized deviation. The primary baseline implementation uses Food
Compass 2.0 scores, and the primary deviation cap is 12 points. Sensitivity
analyses use 8-point and 15-point caps.

\subsection*{Food Compass 2.0 baseline prior}

Food Compass 2.0 scores were extracted from Supplementary Table S5 of the
Food Compass 2.0 article and supplement \cite{barrett2024foodcompass2}. The table contains
9,273 foods and beverages consumed by US adults in NHANES 2001-2002 through
2017-2018. We parsed the supplementary PDF into a machine-readable table and
recorded an extraction audit containing row count, score range checks,
difference checks, duplicate food codes, missing comparison labels and malformed
rows. The parser preserves all extracted rows and treats duplicate food codes as
join risks rather than resolving them silently.

\subsection*{FNDDS nutrient-vector alignment}

Nutrient vectors are matched to the historical food universe used by Food
Compass 2.0. We anchored Food Compass 2.0 scores to the 9,273 foods and
beverages reported in NHANES 2001-2002 through 2017-2018 and matched nutrient
vectors using cycle-specific USDA FNDDS releases from 2001-2002 to 2017-2018,
preserving FNDDS food codes, descriptions and release years to reduce
cross-cycle food-code drift. Current FoodData Central Survey Foods, Foundation
Foods or Branded Foods releases are not used to replace these historical FNDDS
cycles in the primary analysis. They may be used only for future extension or
sensitivity analyses with separate provenance records.

\subsection*{Personalized microbiome calibration}

For each individual, microbiome-derived nutrient weights \(\beta_i\) are applied
to standardized food nutrient vectors \(X_j\) to obtain a raw personalized
response \(R_{ij} = \beta_i X_j\). The response is decomposed into
MAC/SCFA/plant-matrix, lipid/bile-acid/TMAO and residual channels. The primary
reported channel attributions are the expert-revised MAC and lipid channels.
Carbohydrate is reserved for sensitivity analysis, zinc and copper are excluded
from the primary microbiome masks, vitamin A RAE is excluded and retinol is
retained as the lipid-channel marker.

The raw response is robustly standardized across individuals for each food and
mapped to a bounded deviation:
\[
  D_{ij} = 12 \tanh(Z_{ij}/2).
\]
The deviation is centered within each food before final clipping to the 1-100
score range. This centering forces the population mean of the personalized
deviation toward zero, so the calibrated score remains anchored to the
population NPS prior.

\subsection*{Synthetic digital-gut-twin benchmark}

We built a synthetic digital-gut-twin benchmark to test the scoring architecture
under known ground truth. The simulator generated individual MAC and lipid
capacities, real-valued food substrate exposures, a universal food-quality
prior and microbiome-conditioned individual response terms. The known response
combined universal food quality, personalized MAC and lipid effects and random
noise. We compared FCS2-only scoring, unanchored microbiome scoring, anchored
GMNPS, random microbiome weights, shuffled microbiome weights, the original
mask and the expert-revised mask. The benchmark tested whether anchored GMNPS
could improve recovery of individual response while preserving the universal
NPS ranking better than an unanchored microbiome score.

\subsection*{Retrospective external biological anchor}

We used the published ZOE Microbiome Health Ranking 2025 as an external
biological anchor for the microbiome calibration layer
\cite{asnicar2025gut}. Publicly available species-level MetaPhlAn profiles and
age-, sex- and BMI-limited metadata can support retrospective plausibility
analyses. Restricted individual-level ZOE clinical, dietary and intervention
data were not used in the current computational benchmark; therefore, these
analyses were not designed to estimate clinical efficacy. Microbial features
should be summarized a priori using published favourable and unfavourable ranks
to avoid outcome-driven feature selection.
